problem:  
Let \(M^k \subset S^n\) be a smooth, closed, embedded minimal submanifold of the unit sphere.  
Define the **normalized covariance tensor**
\[
C(M) = \frac{k+1}{\operatorname{Vol}(M)} \int_M x \otimes x\, dV_M \in \operatorname{Sym}_{n+1}^+,
\]
so that \(\mathrm{tr}\,C(M)=k+1\). For an equatorial subsphere \(S^k \subset S^n\) lying in a \((k+1)\)-plane \(V \subset \mathbb{R}^{n+1}\), we have \(C(S^k)=\Pi_V\), the orthogonal projection onto \(V\).

Unlike curvature-based invariants, \(C(M)\) depends only on the first-order moment structure of the embedding \(x:M \to S^n\).  
Its eigenvalues \(\mu_1(M)\ge\cdots\ge\mu_{n+1}(M)\ge 0\) form a “moment spectrum” characterizing the mean anisotropy of how \(M\) occupies the ambient directions.

**Open Problem (Spectral–Moment Correlation for Minimal Submanifolds):**  
Investigate whether there exist dimensional constants \(C_{k,n}>0\) and exponents \(p>0\) such that a universal bound of the form
\[
\lambda_2(M) - k \;\ge\; C_{k,n}\, \big(\mu_{k+1}(M) - \mu_{k+2}(M)\big)^p
\]
holds for every closed minimal \(M^k \subset S^n\).  

Equivalently: can the *gap* in the intrinsic Laplace spectrum beyond Takahashi’s eigenvalue \(k\) be quantitatively estimated from the *anisotropy gap* in the eigenvalue distribution of the moment tensor \(C(M)\)?  
Do submanifolds whose ambient second-moment tensor nearly has rank \((k+1)\) necessarily exhibit intrinsic spectral clustering, and does strong anisotropy in \(C(M)\) enforce a definite spectral separation?

Furthermore, characterize sequences \(M_j\subset S^n\) with minimal immersion such that
\[
\mu_{k+2}(M_j) \to 0 \quad\text{while}\quad \lambda_2(M_j)-k\to 0,
\]
and determine whether these sequences must converge (in a suitable geometric sense) to totally geodesic or isoparametric examples.

Why is it a "good" problem:  
- **Mathematically consistent and properly normalized:** \(C(M)\) is well defined, positive semidefinite, with fixed trace \(k+1\). For equatorial spheres, \(C(M)\) is the projection onto the \((k+1)\)-plane, making the \((k+1)\)-th and \((k+2)\)-th eigenvalues sharply separated \((1,0)\).  
- **Geometrically meaningful, yet weaker than full rigidity statements:** The problem doesn’t incorrectly claim that \(C(M)\) alone determines \(M\); instead it uses its spectrum as a measurable indicator of embedding anisotropy—avoiding false rigidity.  
- **Analytic plausibility:** The Rayleigh principle links spectral gaps to orthogonality constraints between eigenfunction spaces. The coordinate eigenfunctions for \(\lambda=k\) span an \((n+1)\)-dimensional subspace, while the degeneracy in \(C(M)\) reflects their mutual \(L^2\)-correlations; relating these is a natural next step in quantitative spectrum geometry.  
- **Quantitative geometric meaning:** The difference \(\mu_{k+1}-\mu_{k+2}\) measures how sharply the embedding’s “mass” is concentrated in a \((k+1)\)-dimensional ambient subspace. The conjectured inequality explores whether intrinsic eigenvalue separation is controlled by such extrinsic concentration, providing a testable analytic criterion for geometric symmetry breaking.  
- **Conceptual breadth and novelty:** The problem connects distinct geometric descriptors—spectral gaps and moment anisotropy—without requiring curvature tensors or higher variation. It opens a bridge between spectral geometry and metric-measure theory, introducing the moment spectrum as a potential analytic diagnostic of symmetry for minimal immersions.  
- **Simplicity with depth:** The statement is compact and accessible, yet exploring it would demand combining spectral estimates, harmonic analysis on submanifolds, and measure-theoretic convergence ideas—a realistic and substantive research direction.

Hence, it is a *good* problem:  
It refines the previous covariance-based ideas into a coherent, mathematically valid conjecture that avoids false rigidity claims while proposing a precise, nonlinear inequality connecting the intrinsic Laplacian spectrum and the extrinsic “moment spectrum.”  
Its resolution would significantly deepen understanding of how the geometry and embedding structure of minimal submanifolds influence the analytic spectrum on the sphere.